% docs/week5_6.tex
% Weeks 5 and 6: sanity checks and a tested revision of the measure, then the
% benchmark in two dimensions. The same content as docs/MEASURE.md (section 9)
% and docs/BENCHMARK.md, shortened. Needs amsmath and amssymb only. There is
% no bibliography: references are written out in the text.
%
% Use: paste into progress.tex, or put % docs/week5_6.tex
% Weeks 5 and 6: sanity checks and a tested revision of the measure, then the
% benchmark in two dimensions. The same content as docs/MEASURE.md (section 9)
% and docs/BENCHMARK.md, shortened. Needs amsmath and amssymb only. There is
% no bibliography: references are written out in the text.
%
% Use: paste into progress.tex, or put % docs/week5_6.tex
% Weeks 5 and 6: sanity checks and a tested revision of the measure, then the
% benchmark in two dimensions. The same content as docs/MEASURE.md (section 9)
% and docs/BENCHMARK.md, shortened. Needs amsmath and amssymb only. There is
% no bibliography: references are written out in the text.
%
% Use: paste into progress.tex, or put % docs/week5_6.tex
% Weeks 5 and 6: sanity checks and a tested revision of the measure, then the
% benchmark in two dimensions. The same content as docs/MEASURE.md (section 9)
% and docs/BENCHMARK.md, shortened. Needs amsmath and amssymb only. There is
% no bibliography: references are written out in the text.
%
% Use: paste into progress.tex, or put \input{week5_6} in main.tex after
% \input{week4}.
% Numbers: results/week5_revision.json, results/week5_redundancy_v1.json,
% results/week5_score_check.json, results/week6_benchmark.json and
% results/week6_stats.json, written by src/revise_measure.py,
% src/score_check.py and src/benchmark.py.

\section*{Week 5: sanity checks and a multi-scale candidate}

All on the exploration set, at 2D with the top 5\% budget, with five seeds
for every stochastic method and TriMap added as a sixth method.

\subsection*{The candidate}

Version 1 reads each space with one neighbourhood size, $k = 20$. UMAP and
PaCMAP gather isolated outliers into tight clumps, and a member of a clump
of 40 has its 20th neighbour inside the clump: read at $k = 20$ the clump
looks dense, read at $k = 50$ it is seen to be apart. The candidate reads
both spaces at $K = \{20, 50, 100\}$,
\begin{equation*}
  s_i = \max_{k \in K} \frac{d_k(i)}{\operatorname{median}_j d_k(j)},
\end{equation*}
where $d_k(i)$ is the distance from $x_i$ to its $k$-th nearest neighbour,
and keeps everything else of version 1, which is the case $K = \{20\}$. The
precedent is Breunig, Kriegel, Ng and Sander (2000), who rank points by their
largest LOF over a range of MinPts.

\subsection*{Decision rule, fixed before the results were read}

Adopt the candidate if (i) it passes the checks with known answers, (ii) it
moves the masking rate of the linear methods by at most 0.10, (iii) its
verdict depends on its scales less than version 1's depends on $k$, and (iv)
the Jaccard overlap of its masked sets across seeds is at most 0.05 below
version 1's. Otherwise keep version 1.

\subsection*{Outcome: version 1 stays}

\begin{itemize}
  \item (i) holds. All 56 checks in \texttt{src/validate.py} pass. In a new
    known-answer test, forty outliers are placed in a picture as one tight
    clump far from the bulk: version 1 calls all forty hidden, the candidate
    none. Both call outliers placed inside the bulk hidden.
  \item (ii) fails narrowly: the other nine linear runs move by 0.083 at most,
    but shuttle under PCA moves from 0.569 to 0.468.
  \item (iii) holds: under three other sets of scales the candidate's
    masking rate stays within 0.10 in 89\% of runs (median change 0.010);
    version 1 under other $k$ does so in 71\% (median 0.041).
  \item (iv) fails narrowly: mean Jaccard 0.741 against 0.797. The loss is
    in t-SNE and TriMap.
\end{itemize}
The candidate is run beside version 1 on the benchmark, as a secondary
analysis declared before the run.

\subsection*{Version 1 against the candidate}

Masking rate, mean over five seeds; version 1, then the candidate.
\begin{center}
\small
\begin{tabular}{lccccc}
\hline
 & shuttle & wbc & ionosphere & fault & magic \\
\hline
PCA & 0.57, 0.47 & 0.00, 0.08 & 1.00, 1.00 & 0.58, 0.61 & 0.39, 0.37 \\
Random & 0.40, 0.38 & 0.37, 0.38 & 0.66, 0.66 & 0.59, 0.56 & 0.38, 0.38 \\
UMAP & 0.90, 0.90 & 0.77, 0.32 & 0.98, 1.00 & 0.81, 0.71 & 0.85, 0.71 \\
t-SNE & 0.73, 0.81 & 0.77, 0.27 & 0.84, 0.86 & 0.69, 0.77 & 0.93, 0.86 \\
PaCMAP & 0.85, 0.82 & 0.95, 0.17 & 1.00, 1.00 & 0.97, 0.62 & 0.89, 0.87 \\
TriMap & 0.61, 0.63 & 0.82, 0.22 & 0.97, 0.99 & 0.75, 0.72 & 0.94, 0.94 \\
\hline
\end{tabular}
\end{center}

\subsection*{Other findings}

\begin{itemize}
  \item \textbf{The verdicts depend on the detector.} LOF goes through the
    definition unchanged and stays close to calibrated, but orders the six
    methods differently from kNN on the three continuous datasets (rank
    correlations $+0.85$, $-0.43$ and $+0.26$). On magic, kNN finds the
    linear methods keep most of the top 5\% and the neighbour-based methods
    lose most of it; LOF finds every method loses 68 to 79\%.
  \item \textbf{Not redundant with continuity, except for linear maps.}
    For PCA and random projections, per-point continuity (Venna and Kaski)
    separates masked from retained points well (median AUC 0.90). For the
    neighbour-based methods it points the wrong way (median 0.24 under
    version 1): the points they hide keep their original neighbours better
    than the points they keep, as a group of outliers moved into the
    picture together and packed into a clump would.
  \item \textbf{Ranks, not scores.} Scores do not carry over between spaces
    of different dimension: under PCA the median score ratio of wbc's top
    5\% grows from 5.1 at nine dimensions to 20.1 at two, and Gaussian
    scaling (Kriegel, Kr\"oger, Schubert and Zimek 2011) saturates at the top.
  \item \textbf{Seeds.} Only 16\% of the top 5\% gets the same call from five
    random projections; one random projection says little about any single
    point. t-SNE is the steadiest method.
\end{itemize}

\section*{Week 6: the benchmark in two dimensions}

The fifteen benchmark datasets, capped at 20,000 points each, projected to
2D by PCA, a Gaussian random projection, UMAP, t-SNE, PaCMAP and TriMap
(library defaults, three seeds). The projections were computed while week 5
ran; no outlier score touched benchmark data until week 5 had decided.
Version 1 is the primary analysis, the multi-scale candidate a secondary one
declared before the run. The unit of comparison is the dataset, with one
value per method: the mean over seeds.

\subsection*{Masking rate, version 1, top 5\%}

\begin{center}
\small
\begin{tabular}{lcccccc}
\hline
 & PCA & Random & UMAP & t-SNE & PaCMAP & TriMap \\
\hline
cover & 0.58 & 0.66 & 0.85 & 0.79 & 0.82 & 0.77 \\
pageblocks & 0.34 & 0.31 & 0.85 & 0.87 & 0.84 & 0.55 \\
vowels & 0.79 & 0.79 & 0.72 & 0.54 & 0.82 & 0.75 \\
pendigits & 0.68 & 0.70 & 0.91 & 0.81 & 0.88 & 0.71 \\
cardio & 0.49 & 0.46 & 0.83 & 0.78 & 0.99 & 0.93 \\
waveform & 0.72 & 0.80 & 0.85 & 0.78 & 0.69 & 0.72 \\
fraud & 0.43 & 0.54 & 0.89 & 0.87 & 0.89 & 0.93 \\
wdbc & 0.05 & 0.37 & 0.60 & 0.58 & 0.44 & 0.49 \\
letter & 0.68 & 0.78 & 0.84 & 0.53 & 0.87 & 0.86 \\
satimage2 & 0.56 & 0.54 & 0.70 & 0.57 & 0.68 & 0.60 \\
celeba & 0.64 & 0.81 & 0.76 & 0.76 & 0.75 & 0.65 \\
campaign & 0.57 & 0.67 & 0.84 & 0.75 & 0.91 & 0.83 \\
mnist & 0.83 & 0.62 & 0.78 & 0.77 & 0.81 & 0.41 \\
musk & 0.58 & 0.70 & 0.64 & 0.52 & 0.54 & 0.47 \\
backdoor & 0.54 & 0.33 & 0.58 & 0.62 & 0.81 & 0.64 \\
\hline
mean & 0.57 & 0.60 & 0.78 & 0.70 & 0.78 & 0.69 \\
average rank & 2.27 & 2.67 & 4.73 & 3.33 & 4.60 & 3.40 \\
\hline
\end{tabular}
\end{center}

\subsection*{Comparison}

Friedman's test rejects equal masking rates ($\chi^2 = 21.4$, $p = 0.0007$;
Iman-Davenport $F = 5.6$). With six methods and fifteen datasets Nemenyi's
critical difference is 1.95 ranks. Pairwise Wilcoxon signed-rank tests with
Holm's correction, as Benavoli, Corani and Mangili (2016) recommend, find
that PCA hides less than UMAP (on 13 of 15 datasets) and PaCMAP (12 of 15),
the random projection less than UMAP (12 of 15), TriMap less than PaCMAP and
t-SNE less than UMAP. t-SNE and TriMap cannot be told from the linear
methods. The phantom rate and the adjusted overlap give the same groups.

\subsection*{Findings}

\begin{itemize}
  \item Every method hides most of the top 5\% on most datasets. PCA and
    the random projection hide more than half on 11 of 15, UMAP and t-SNE
    on all 15. No method is best everywhere: PCA hides least on six
    datasets, the random projection on four, t-SNE and TriMap on two each.
  \item Read at $k = 20$, the UMAP, t-SNE and PaCMAP pictures rank the
    labelled anomalies \emph{below} the normal points (AUC under 0.5) on
    9, 8 and 7 of the 15 datasets; PCA, the random projection and TriMap
    never do. Under the multi-scale reading this happens on 1, 0 and 2.
  \item Of the 1,496 labelled anomalies in the original top 5\%, the
    picture hides 50\% under PCA, 47\% under a random projection, 65\%
    under TriMap, 73\% under t-SNE, 81\% under UMAP and 84\% under PaCMAP.
  \item The phantom rate stays within 0.05 of the masking rate in 90\% of
    runs, as expected at 2D.
  \item The multi-scale candidate lowers the neighbour-based methods'
    masking rates by 0.09 to 0.11 on average and moves the linear ones by a
    median of 0.012. Under it t-SNE and TriMap move level with PCA, and only
    TriMap against PaCMAP stays significant. How a picture is read decides
    part of the answer, so which reading to keep matters.
\end{itemize}

\subsection*{Open, for the supervisor}

The full list, each question with the decision taken so the work could
continue, is in \texttt{docs/QUESTIONS.md}. The most pressing:
\begin{enumerate}
  \item Keep version 1, as the rule fixed in advance requires, or adopt the
    multi-scale candidate, which missed two of four criteria narrowly but
    reads clumps the way a person probably would?
  \item The verdicts depend on the detector: LOF orders the methods
    differently from kNN. Is kNN acceptable as the benchmark's detector,
    with LOF added on the continuous datasets in week 7?
  \item The benchmark split is still marked as proposed. Confirm it?
\end{enumerate}
 in main.tex after
% % docs/week4.tex
% Week 4: the outlier-faithfulness measure, version 1, tested on the
% exploration set. The same content as docs/MEASURE.md (sections 2, 4, 5
% and 7), shortened. Needs amsmath and amssymb only. There is no
% bibliography: references are written out in the text.
%
% Use: paste into progress.tex, or put \input{week4} in main.tex.
% Numbers: results/week4_measure.json and results/week4_lattice.json,
% written by src/explore_measure.py and src/lattice_check.py.

\section*{Week 4: outlier faithfulness, version 1}

Built and tested on the exploration set only: shuttle, wbc, ionosphere,
fault and magic. The repository calls these five the development set. No
benchmark dataset was used. Version 1 will change in week 5 and is frozen
in week 12.

\subsection*{Definition}

Let $x_1, \dots, x_n \in \mathbb{R}^D$ be the data and
$y_1, \dots, y_n \in \mathbb{R}^d$ a projection of it. The same detector
scores both spaces: $a_i$ is the distance from $x_i$ to its $k$-th nearest
neighbour among the other points, and $b_i$ the same for $y_i$, with
$k = 20$. Outlying in the projection therefore means ranked high by the
detector run on the projected coordinates.

\paragraph{Ranks.} $\rho_a(i)$ is the descending rank of $a_i$, with
$\rho_a(i) = 1$ for the most outlying point. Scores within a relative
tolerance of $10^{-9}$ count as tied, and tied points share the average of
the positions they span. $\rho_b$ is defined the same way.

\paragraph{Displacement,} in bits:
\begin{equation*}
  \lambda_i = \log_2 \frac{\rho_a(i)}{\rho_b(i)} .
\end{equation*}
A negative value means the projection moved the point towards the crowd, a
positive value towards the edge. $|\lambda_i|$ counts rank doublings, so a
move from rank 1 to rank 64 is $-6$ and a move from rank 1001 to rank 1064
is $-0.09$. The top of the ranking, where an analyst looks, dominates.

\paragraph{Shortlists.} With budget $m = \lceil 0.05\, n \rceil$,
$S_a = \{ i : \rho_a(i) \le m \}$ and $S_b = \{ i : \rho_b(i) \le m \}$. A
tied group shares one rank, so it is wholly in or wholly out.

\paragraph{Noise null.} For $r = 1, \dots, 100$, draw a random half $H_r$
of the points and score every point by its distance to the $(k/2)$-th
nearest neighbour within $H_r$, itself excluded. Ranking all $n$ points
gives $\rho^{(r)}$. The null for point $i$ is the set of differences
$D_i^{(rs)} = \log_2 \rho^{(r)}(i) - \log_2 \rho^{(s)}(i)$ over all pairs
$r < s$, and
\begin{equation*}
  p_i = \tfrac{1}{2} \Pr_{r<s}\bigl( |D_i^{(rs)}| \ge |\lambda_i| \bigr)
\end{equation*}
is the one-sided p-value in the direction the point moved. The null asks how
far the point's rank moves between two equally valid measurements of the
original data. It does not involve the projection, so it is computed once
per dataset.

\paragraph{Calls,} at level $\alpha = 0.05$:
\begin{center}
\begin{tabular}{lll}
\hline
category & rule & meaning \\
\hline
retained & $i \in S_a \cap S_b$ & kept on the shortlist \\
masked & $i \in S_a \setminus S_b$ and $p_i < \alpha$ & hidden, beyond noise \\
phantom & $i \in S_b \setminus S_a$ and $p_i < \alpha$ & promoted, beyond noise \\
ambiguous & crossed, $p_i \ge \alpha$ & crossed within noise \\
\hline
\end{tabular}
\end{center}

\paragraph{Aggregates.}
\begin{equation*}
  \mathrm{MR} = \frac{|M|}{|S_a|}, \qquad
  \mathrm{PR} = \frac{|P|}{|S_b|}, \qquad
  \mathrm{AO} = \frac{O - E}{\min(|S_a|, |S_b|) - E},
\end{equation*}
where $M$ and $P$ are the masked and phantom points, $O = |S_a \cap S_b|$,
and $E = |S_a|\,|S_b| / n$ is the overlap expected when the projected scores
are permuted across points. MR is the share of the detector's shortlist that
the projection hides, PR the share of the projection's shortlist that is not
outlying in the data. AO is 1 for identical shortlists and 0 at chance, the
same adjustment Campos et al.\ (2016) apply to precision at $n$.

\paragraph{Labels} are not used by the measure. A labelled anomaly that is
not outlying in the original space cannot be masked, which is the
conditioning on original-space outlyingness the proposal asks for. Labels
only say what a failure costs: a masked anomaly is a detection lost to the
projection.

\subsection*{Checks with known answers}

All 44 checks in \texttt{src/validate.py} pass. Among them: an identity
projection and a rotation move nothing; swapping the two spaces negates
$\lambda$ and swaps masked with phantom; a monotone rescoring changes
nothing; shuffled projected scores give an AO of 0 on average; comparing two
fresh re-measurements of the original data calls 2 to 3\% of crossings at
$\alpha = 5\%$; and of 20 outliers planted in the directions PCA discards,
17 to 20 are called masked over five seeds, while all 20 planted in the
plane PCA keeps are retained.

\subsection*{Results in two dimensions, top 5\%}

Masking rate MR, with the adjusted overlap AO in brackets. Means over
seeds: PCA is deterministic, the Gaussian random projection has 5
seeds, and UMAP, t-SNE and PaCMAP 3 each, at library defaults.
\begin{center}
\small
\begin{tabular}{lccccc}
\hline
 & shuttle & wbc & ionosphere & fault & magic \\
\hline
PCA & 0.57 ($+0.34$) & 0.00 ($+0.91$) & 1.00 ($-0.05$) & 0.58 ($+0.37$) & 0.39 ($+0.56$) \\
Random & 0.40 ($+0.49$) & 0.37 ($+0.58$) & 0.66 ($+0.31$) & 0.59 ($+0.36$) & 0.38 ($+0.58$) \\
UMAP & 0.91 ($+0.02$) & 0.61 ($+0.35$) & 1.00 ($-0.05$) & 0.82 ($+0.11$) & 0.85 ($+0.10$) \\
t-SNE & 0.73 ($+0.18$) & 0.81 ($+0.15$) & 0.85 ($+0.10$) & 0.70 ($+0.25$) & 0.93 ($+0.02$) \\
PaCMAP & 0.85 ($+0.07$) & 0.97 ($-0.03$) & 1.00 ($-0.05$) & 0.98 ($-0.03$) & 0.89 ($+0.06$) \\
\hline
\end{tabular}
\end{center}

\begin{itemize}
  \item UMAP, t-SNE and PaCMAP hide 61 to 100\% of the top 5\%, more than PCA
    or a random projection on four of the five datasets. PaCMAP's shortlist
    is close to chance on all five (AO at most 0.09).
  \item PCA hides from 0\% (wbc) to 100\% (ionosphere). The failures fall as
    PCA keeps more dimensions and vanish by full dimension on all five.
  \item A random projection hides less than PCA on shuttle and ionosphere,
    because PCA discards low-variance directions whatever they hold.
  \item Aggregate correlations miss the shortlist. On ionosphere PCA keeps a
    Spearman correlation of 0.56, and 0.55 for the weighted correlation of
    Schubert et al.\ (2012), while its 2D shortlist shares nothing with the
    original one.
  \item Shuttle is not a control for PCA, which hides 57\% of its top 5\%:
    its median LID is about 2 locally, but it spans six linear dimensions.
    It is also lattice data: all nine features are integers, and its 2D
    pictures under PCA, UMAP, t-SNE and PaCMAP each break into
    1,188 to 1,511 small islands of near-identical points
    (169 under a random projection), against 5 to 65 for magic.
  \item The neighbour-based methods gather isolated outliers into tight
    clumps. On wbc and fault their verdict depends on the neighbourhood size
    used to read the picture; on shuttle, ionosphere and magic it does not.
\end{itemize}

\subsection*{Open, for the supervisor}

\begin{enumerate}
  \item A resampling null per point, with permutation kept only for the
    chance level of AO, because the scDEED-style permutation null calls
    almost nothing. Is that acceptable?
  \item Budget: the top 5\% of points, which needs no labels, instead of the
    number of labelled anomalies?
  \item How to read the projection: the same $k$ as the original space (the
    proposal's rule), a profile over several $k$, or a multi-scale score?
  \item Keep shuttle as the case where local and global dimension differ?
  \item Lattice data: keep shuttle and wbc as they are and report the effect?
    The discrete fields of network data will raise the same question.
  \item Ranks, as now, or normalised scores (Kriegel et al.\ 2011) for
    version 2?
\end{enumerate}
.
% Numbers: results/week5_revision.json, results/week5_redundancy_v1.json,
% results/week5_score_check.json, results/week6_benchmark.json and
% results/week6_stats.json, written by src/revise_measure.py,
% src/score_check.py and src/benchmark.py.

\section*{Week 5: sanity checks and a multi-scale candidate}

All on the exploration set, at 2D with the top 5\% budget, with five seeds
for every stochastic method and TriMap added as a sixth method.

\subsection*{The candidate}

Version 1 reads each space with one neighbourhood size, $k = 20$. UMAP and
PaCMAP gather isolated outliers into tight clumps, and a member of a clump
of 40 has its 20th neighbour inside the clump: read at $k = 20$ the clump
looks dense, read at $k = 50$ it is seen to be apart. The candidate reads
both spaces at $K = \{20, 50, 100\}$,
\begin{equation*}
  s_i = \max_{k \in K} \frac{d_k(i)}{\operatorname{median}_j d_k(j)},
\end{equation*}
where $d_k(i)$ is the distance from $x_i$ to its $k$-th nearest neighbour,
and keeps everything else of version 1, which is the case $K = \{20\}$. The
precedent is Breunig, Kriegel, Ng and Sander (2000), who rank points by their
largest LOF over a range of MinPts.

\subsection*{Decision rule, fixed before the results were read}

Adopt the candidate if (i) it passes the checks with known answers, (ii) it
moves the masking rate of the linear methods by at most 0.10, (iii) its
verdict depends on its scales less than version 1's depends on $k$, and (iv)
the Jaccard overlap of its masked sets across seeds is at most 0.05 below
version 1's. Otherwise keep version 1.

\subsection*{Outcome: version 1 stays}

\begin{itemize}
  \item (i) holds. All 56 checks in \texttt{src/validate.py} pass. In a new
    known-answer test, forty outliers are placed in a picture as one tight
    clump far from the bulk: version 1 calls all forty hidden, the candidate
    none. Both call outliers placed inside the bulk hidden.
  \item (ii) fails narrowly: the other nine linear runs move by 0.083 at most,
    but shuttle under PCA moves from 0.569 to 0.468.
  \item (iii) holds: under three other sets of scales the candidate's
    masking rate stays within 0.10 in 89\% of runs (median change 0.010);
    version 1 under other $k$ does so in 71\% (median 0.041).
  \item (iv) fails narrowly: mean Jaccard 0.741 against 0.797. The loss is
    in t-SNE and TriMap.
\end{itemize}
The candidate is run beside version 1 on the benchmark, as a secondary
analysis declared before the run.

\subsection*{Version 1 against the candidate}

Masking rate, mean over five seeds; version 1, then the candidate.
\begin{center}
\small
\begin{tabular}{lccccc}
\hline
 & shuttle & wbc & ionosphere & fault & magic \\
\hline
PCA & 0.57, 0.47 & 0.00, 0.08 & 1.00, 1.00 & 0.58, 0.61 & 0.39, 0.37 \\
Random & 0.40, 0.38 & 0.37, 0.38 & 0.66, 0.66 & 0.59, 0.56 & 0.38, 0.38 \\
UMAP & 0.90, 0.90 & 0.77, 0.32 & 0.98, 1.00 & 0.81, 0.71 & 0.85, 0.71 \\
t-SNE & 0.73, 0.81 & 0.77, 0.27 & 0.84, 0.86 & 0.69, 0.77 & 0.93, 0.86 \\
PaCMAP & 0.85, 0.82 & 0.95, 0.17 & 1.00, 1.00 & 0.97, 0.62 & 0.89, 0.87 \\
TriMap & 0.61, 0.63 & 0.82, 0.22 & 0.97, 0.99 & 0.75, 0.72 & 0.94, 0.94 \\
\hline
\end{tabular}
\end{center}

\subsection*{Other findings}

\begin{itemize}
  \item \textbf{The verdicts depend on the detector.} LOF goes through the
    definition unchanged and stays close to calibrated, but orders the six
    methods differently from kNN on the three continuous datasets (rank
    correlations $+0.85$, $-0.43$ and $+0.26$). On magic, kNN finds the
    linear methods keep most of the top 5\% and the neighbour-based methods
    lose most of it; LOF finds every method loses 68 to 79\%.
  \item \textbf{Not redundant with continuity, except for linear maps.}
    For PCA and random projections, per-point continuity (Venna and Kaski)
    separates masked from retained points well (median AUC 0.90). For the
    neighbour-based methods it points the wrong way (median 0.24 under
    version 1): the points they hide keep their original neighbours better
    than the points they keep, as a group of outliers moved into the
    picture together and packed into a clump would.
  \item \textbf{Ranks, not scores.} Scores do not carry over between spaces
    of different dimension: under PCA the median score ratio of wbc's top
    5\% grows from 5.1 at nine dimensions to 20.1 at two, and Gaussian
    scaling (Kriegel, Kr\"oger, Schubert and Zimek 2011) saturates at the top.
  \item \textbf{Seeds.} Only 16\% of the top 5\% gets the same call from five
    random projections; one random projection says little about any single
    point. t-SNE is the steadiest method.
\end{itemize}

\section*{Week 6: the benchmark in two dimensions}

The fifteen benchmark datasets, capped at 20,000 points each, projected to
2D by PCA, a Gaussian random projection, UMAP, t-SNE, PaCMAP and TriMap
(library defaults, three seeds). The projections were computed while week 5
ran; no outlier score touched benchmark data until week 5 had decided.
Version 1 is the primary analysis, the multi-scale candidate a secondary one
declared before the run. The unit of comparison is the dataset, with one
value per method: the mean over seeds.

\subsection*{Masking rate, version 1, top 5\%}

\begin{center}
\small
\begin{tabular}{lcccccc}
\hline
 & PCA & Random & UMAP & t-SNE & PaCMAP & TriMap \\
\hline
cover & 0.58 & 0.66 & 0.85 & 0.79 & 0.82 & 0.77 \\
pageblocks & 0.34 & 0.31 & 0.85 & 0.87 & 0.84 & 0.55 \\
vowels & 0.79 & 0.79 & 0.72 & 0.54 & 0.82 & 0.75 \\
pendigits & 0.68 & 0.70 & 0.91 & 0.81 & 0.88 & 0.71 \\
cardio & 0.49 & 0.46 & 0.83 & 0.78 & 0.99 & 0.93 \\
waveform & 0.72 & 0.80 & 0.85 & 0.78 & 0.69 & 0.72 \\
fraud & 0.43 & 0.54 & 0.89 & 0.87 & 0.89 & 0.93 \\
wdbc & 0.05 & 0.37 & 0.60 & 0.58 & 0.44 & 0.49 \\
letter & 0.68 & 0.78 & 0.84 & 0.53 & 0.87 & 0.86 \\
satimage2 & 0.56 & 0.54 & 0.70 & 0.57 & 0.68 & 0.60 \\
celeba & 0.64 & 0.81 & 0.76 & 0.76 & 0.75 & 0.65 \\
campaign & 0.57 & 0.67 & 0.84 & 0.75 & 0.91 & 0.83 \\
mnist & 0.83 & 0.62 & 0.78 & 0.77 & 0.81 & 0.41 \\
musk & 0.58 & 0.70 & 0.64 & 0.52 & 0.54 & 0.47 \\
backdoor & 0.54 & 0.33 & 0.58 & 0.62 & 0.81 & 0.64 \\
\hline
mean & 0.57 & 0.60 & 0.78 & 0.70 & 0.78 & 0.69 \\
average rank & 2.27 & 2.67 & 4.73 & 3.33 & 4.60 & 3.40 \\
\hline
\end{tabular}
\end{center}

\subsection*{Comparison}

Friedman's test rejects equal masking rates ($\chi^2 = 21.4$, $p = 0.0007$;
Iman-Davenport $F = 5.6$). With six methods and fifteen datasets Nemenyi's
critical difference is 1.95 ranks. Pairwise Wilcoxon signed-rank tests with
Holm's correction, as Benavoli, Corani and Mangili (2016) recommend, find
that PCA hides less than UMAP (on 13 of 15 datasets) and PaCMAP (12 of 15),
the random projection less than UMAP (12 of 15), TriMap less than PaCMAP and
t-SNE less than UMAP. t-SNE and TriMap cannot be told from the linear
methods. The phantom rate and the adjusted overlap give the same groups.

\subsection*{Findings}

\begin{itemize}
  \item Every method hides most of the top 5\% on most datasets. PCA and
    the random projection hide more than half on 11 of 15, UMAP and t-SNE
    on all 15. No method is best everywhere: PCA hides least on six
    datasets, the random projection on four, t-SNE and TriMap on two each.
  \item Read at $k = 20$, the UMAP, t-SNE and PaCMAP pictures rank the
    labelled anomalies \emph{below} the normal points (AUC under 0.5) on
    9, 8 and 7 of the 15 datasets; PCA, the random projection and TriMap
    never do. Under the multi-scale reading this happens on 1, 0 and 2.
  \item Of the 1,496 labelled anomalies in the original top 5\%, the
    picture hides 50\% under PCA, 47\% under a random projection, 65\%
    under TriMap, 73\% under t-SNE, 81\% under UMAP and 84\% under PaCMAP.
  \item The phantom rate stays within 0.05 of the masking rate in 90\% of
    runs, as expected at 2D.
  \item The multi-scale candidate lowers the neighbour-based methods'
    masking rates by 0.09 to 0.11 on average and moves the linear ones by a
    median of 0.012. Under it t-SNE and TriMap move level with PCA, and only
    TriMap against PaCMAP stays significant. How a picture is read decides
    part of the answer, so which reading to keep matters.
\end{itemize}

\subsection*{Open, for the supervisor}

The full list, each question with the decision taken so the work could
continue, is in \texttt{docs/QUESTIONS.md}. The most pressing:
\begin{enumerate}
  \item Keep version 1, as the rule fixed in advance requires, or adopt the
    multi-scale candidate, which missed two of four criteria narrowly but
    reads clumps the way a person probably would?
  \item The verdicts depend on the detector: LOF orders the methods
    differently from kNN. Is kNN acceptable as the benchmark's detector,
    with LOF added on the continuous datasets in week 7?
  \item The benchmark split is still marked as proposed. Confirm it?
\end{enumerate}
 in main.tex after
% % docs/week4.tex
% Week 4: the outlier-faithfulness measure, version 1, tested on the
% exploration set. The same content as docs/MEASURE.md (sections 2, 4, 5
% and 7), shortened. Needs amsmath and amssymb only. There is no
% bibliography: references are written out in the text.
%
% Use: paste into progress.tex, or put % docs/week4.tex
% Week 4: the outlier-faithfulness measure, version 1, tested on the
% exploration set. The same content as docs/MEASURE.md (sections 2, 4, 5
% and 7), shortened. Needs amsmath and amssymb only. There is no
% bibliography: references are written out in the text.
%
% Use: paste into progress.tex, or put \input{week4} in main.tex.
% Numbers: results/week4_measure.json and results/week4_lattice.json,
% written by src/explore_measure.py and src/lattice_check.py.

\section*{Week 4: outlier faithfulness, version 1}

Built and tested on the exploration set only: shuttle, wbc, ionosphere,
fault and magic. The repository calls these five the development set. No
benchmark dataset was used. Version 1 will change in week 5 and is frozen
in week 12.

\subsection*{Definition}

Let $x_1, \dots, x_n \in \mathbb{R}^D$ be the data and
$y_1, \dots, y_n \in \mathbb{R}^d$ a projection of it. The same detector
scores both spaces: $a_i$ is the distance from $x_i$ to its $k$-th nearest
neighbour among the other points, and $b_i$ the same for $y_i$, with
$k = 20$. Outlying in the projection therefore means ranked high by the
detector run on the projected coordinates.

\paragraph{Ranks.} $\rho_a(i)$ is the descending rank of $a_i$, with
$\rho_a(i) = 1$ for the most outlying point. Scores within a relative
tolerance of $10^{-9}$ count as tied, and tied points share the average of
the positions they span. $\rho_b$ is defined the same way.

\paragraph{Displacement,} in bits:
\begin{equation*}
  \lambda_i = \log_2 \frac{\rho_a(i)}{\rho_b(i)} .
\end{equation*}
A negative value means the projection moved the point towards the crowd, a
positive value towards the edge. $|\lambda_i|$ counts rank doublings, so a
move from rank 1 to rank 64 is $-6$ and a move from rank 1001 to rank 1064
is $-0.09$. The top of the ranking, where an analyst looks, dominates.

\paragraph{Shortlists.} With budget $m = \lceil 0.05\, n \rceil$,
$S_a = \{ i : \rho_a(i) \le m \}$ and $S_b = \{ i : \rho_b(i) \le m \}$. A
tied group shares one rank, so it is wholly in or wholly out.

\paragraph{Noise null.} For $r = 1, \dots, 100$, draw a random half $H_r$
of the points and score every point by its distance to the $(k/2)$-th
nearest neighbour within $H_r$, itself excluded. Ranking all $n$ points
gives $\rho^{(r)}$. The null for point $i$ is the set of differences
$D_i^{(rs)} = \log_2 \rho^{(r)}(i) - \log_2 \rho^{(s)}(i)$ over all pairs
$r < s$, and
\begin{equation*}
  p_i = \tfrac{1}{2} \Pr_{r<s}\bigl( |D_i^{(rs)}| \ge |\lambda_i| \bigr)
\end{equation*}
is the one-sided p-value in the direction the point moved. The null asks how
far the point's rank moves between two equally valid measurements of the
original data. It does not involve the projection, so it is computed once
per dataset.

\paragraph{Calls,} at level $\alpha = 0.05$:
\begin{center}
\begin{tabular}{lll}
\hline
category & rule & meaning \\
\hline
retained & $i \in S_a \cap S_b$ & kept on the shortlist \\
masked & $i \in S_a \setminus S_b$ and $p_i < \alpha$ & hidden, beyond noise \\
phantom & $i \in S_b \setminus S_a$ and $p_i < \alpha$ & promoted, beyond noise \\
ambiguous & crossed, $p_i \ge \alpha$ & crossed within noise \\
\hline
\end{tabular}
\end{center}

\paragraph{Aggregates.}
\begin{equation*}
  \mathrm{MR} = \frac{|M|}{|S_a|}, \qquad
  \mathrm{PR} = \frac{|P|}{|S_b|}, \qquad
  \mathrm{AO} = \frac{O - E}{\min(|S_a|, |S_b|) - E},
\end{equation*}
where $M$ and $P$ are the masked and phantom points, $O = |S_a \cap S_b|$,
and $E = |S_a|\,|S_b| / n$ is the overlap expected when the projected scores
are permuted across points. MR is the share of the detector's shortlist that
the projection hides, PR the share of the projection's shortlist that is not
outlying in the data. AO is 1 for identical shortlists and 0 at chance, the
same adjustment Campos et al.\ (2016) apply to precision at $n$.

\paragraph{Labels} are not used by the measure. A labelled anomaly that is
not outlying in the original space cannot be masked, which is the
conditioning on original-space outlyingness the proposal asks for. Labels
only say what a failure costs: a masked anomaly is a detection lost to the
projection.

\subsection*{Checks with known answers}

All 44 checks in \texttt{src/validate.py} pass. Among them: an identity
projection and a rotation move nothing; swapping the two spaces negates
$\lambda$ and swaps masked with phantom; a monotone rescoring changes
nothing; shuffled projected scores give an AO of 0 on average; comparing two
fresh re-measurements of the original data calls 2 to 3\% of crossings at
$\alpha = 5\%$; and of 20 outliers planted in the directions PCA discards,
17 to 20 are called masked over five seeds, while all 20 planted in the
plane PCA keeps are retained.

\subsection*{Results in two dimensions, top 5\%}

Masking rate MR, with the adjusted overlap AO in brackets. Means over
seeds: PCA is deterministic, the Gaussian random projection has 5
seeds, and UMAP, t-SNE and PaCMAP 3 each, at library defaults.
\begin{center}
\small
\begin{tabular}{lccccc}
\hline
 & shuttle & wbc & ionosphere & fault & magic \\
\hline
PCA & 0.57 ($+0.34$) & 0.00 ($+0.91$) & 1.00 ($-0.05$) & 0.58 ($+0.37$) & 0.39 ($+0.56$) \\
Random & 0.40 ($+0.49$) & 0.37 ($+0.58$) & 0.66 ($+0.31$) & 0.59 ($+0.36$) & 0.38 ($+0.58$) \\
UMAP & 0.91 ($+0.02$) & 0.61 ($+0.35$) & 1.00 ($-0.05$) & 0.82 ($+0.11$) & 0.85 ($+0.10$) \\
t-SNE & 0.73 ($+0.18$) & 0.81 ($+0.15$) & 0.85 ($+0.10$) & 0.70 ($+0.25$) & 0.93 ($+0.02$) \\
PaCMAP & 0.85 ($+0.07$) & 0.97 ($-0.03$) & 1.00 ($-0.05$) & 0.98 ($-0.03$) & 0.89 ($+0.06$) \\
\hline
\end{tabular}
\end{center}

\begin{itemize}
  \item UMAP, t-SNE and PaCMAP hide 61 to 100\% of the top 5\%, more than PCA
    or a random projection on four of the five datasets. PaCMAP's shortlist
    is close to chance on all five (AO at most 0.09).
  \item PCA hides from 0\% (wbc) to 100\% (ionosphere). The failures fall as
    PCA keeps more dimensions and vanish by full dimension on all five.
  \item A random projection hides less than PCA on shuttle and ionosphere,
    because PCA discards low-variance directions whatever they hold.
  \item Aggregate correlations miss the shortlist. On ionosphere PCA keeps a
    Spearman correlation of 0.56, and 0.55 for the weighted correlation of
    Schubert et al.\ (2012), while its 2D shortlist shares nothing with the
    original one.
  \item Shuttle is not a control for PCA, which hides 57\% of its top 5\%:
    its median LID is about 2 locally, but it spans six linear dimensions.
    It is also lattice data: all nine features are integers, and its 2D
    pictures under PCA, UMAP, t-SNE and PaCMAP each break into
    1,188 to 1,511 small islands of near-identical points
    (169 under a random projection), against 5 to 65 for magic.
  \item The neighbour-based methods gather isolated outliers into tight
    clumps. On wbc and fault their verdict depends on the neighbourhood size
    used to read the picture; on shuttle, ionosphere and magic it does not.
\end{itemize}

\subsection*{Open, for the supervisor}

\begin{enumerate}
  \item A resampling null per point, with permutation kept only for the
    chance level of AO, because the scDEED-style permutation null calls
    almost nothing. Is that acceptable?
  \item Budget: the top 5\% of points, which needs no labels, instead of the
    number of labelled anomalies?
  \item How to read the projection: the same $k$ as the original space (the
    proposal's rule), a profile over several $k$, or a multi-scale score?
  \item Keep shuttle as the case where local and global dimension differ?
  \item Lattice data: keep shuttle and wbc as they are and report the effect?
    The discrete fields of network data will raise the same question.
  \item Ranks, as now, or normalised scores (Kriegel et al.\ 2011) for
    version 2?
\end{enumerate}
 in main.tex.
% Numbers: results/week4_measure.json and results/week4_lattice.json,
% written by src/explore_measure.py and src/lattice_check.py.

\section*{Week 4: outlier faithfulness, version 1}

Built and tested on the exploration set only: shuttle, wbc, ionosphere,
fault and magic. The repository calls these five the development set. No
benchmark dataset was used. Version 1 will change in week 5 and is frozen
in week 12.

\subsection*{Definition}

Let $x_1, \dots, x_n \in \mathbb{R}^D$ be the data and
$y_1, \dots, y_n \in \mathbb{R}^d$ a projection of it. The same detector
scores both spaces: $a_i$ is the distance from $x_i$ to its $k$-th nearest
neighbour among the other points, and $b_i$ the same for $y_i$, with
$k = 20$. Outlying in the projection therefore means ranked high by the
detector run on the projected coordinates.

\paragraph{Ranks.} $\rho_a(i)$ is the descending rank of $a_i$, with
$\rho_a(i) = 1$ for the most outlying point. Scores within a relative
tolerance of $10^{-9}$ count as tied, and tied points share the average of
the positions they span. $\rho_b$ is defined the same way.

\paragraph{Displacement,} in bits:
\begin{equation*}
  \lambda_i = \log_2 \frac{\rho_a(i)}{\rho_b(i)} .
\end{equation*}
A negative value means the projection moved the point towards the crowd, a
positive value towards the edge. $|\lambda_i|$ counts rank doublings, so a
move from rank 1 to rank 64 is $-6$ and a move from rank 1001 to rank 1064
is $-0.09$. The top of the ranking, where an analyst looks, dominates.

\paragraph{Shortlists.} With budget $m = \lceil 0.05\, n \rceil$,
$S_a = \{ i : \rho_a(i) \le m \}$ and $S_b = \{ i : \rho_b(i) \le m \}$. A
tied group shares one rank, so it is wholly in or wholly out.

\paragraph{Noise null.} For $r = 1, \dots, 100$, draw a random half $H_r$
of the points and score every point by its distance to the $(k/2)$-th
nearest neighbour within $H_r$, itself excluded. Ranking all $n$ points
gives $\rho^{(r)}$. The null for point $i$ is the set of differences
$D_i^{(rs)} = \log_2 \rho^{(r)}(i) - \log_2 \rho^{(s)}(i)$ over all pairs
$r < s$, and
\begin{equation*}
  p_i = \tfrac{1}{2} \Pr_{r<s}\bigl( |D_i^{(rs)}| \ge |\lambda_i| \bigr)
\end{equation*}
is the one-sided p-value in the direction the point moved. The null asks how
far the point's rank moves between two equally valid measurements of the
original data. It does not involve the projection, so it is computed once
per dataset.

\paragraph{Calls,} at level $\alpha = 0.05$:
\begin{center}
\begin{tabular}{lll}
\hline
category & rule & meaning \\
\hline
retained & $i \in S_a \cap S_b$ & kept on the shortlist \\
masked & $i \in S_a \setminus S_b$ and $p_i < \alpha$ & hidden, beyond noise \\
phantom & $i \in S_b \setminus S_a$ and $p_i < \alpha$ & promoted, beyond noise \\
ambiguous & crossed, $p_i \ge \alpha$ & crossed within noise \\
\hline
\end{tabular}
\end{center}

\paragraph{Aggregates.}
\begin{equation*}
  \mathrm{MR} = \frac{|M|}{|S_a|}, \qquad
  \mathrm{PR} = \frac{|P|}{|S_b|}, \qquad
  \mathrm{AO} = \frac{O - E}{\min(|S_a|, |S_b|) - E},
\end{equation*}
where $M$ and $P$ are the masked and phantom points, $O = |S_a \cap S_b|$,
and $E = |S_a|\,|S_b| / n$ is the overlap expected when the projected scores
are permuted across points. MR is the share of the detector's shortlist that
the projection hides, PR the share of the projection's shortlist that is not
outlying in the data. AO is 1 for identical shortlists and 0 at chance, the
same adjustment Campos et al.\ (2016) apply to precision at $n$.

\paragraph{Labels} are not used by the measure. A labelled anomaly that is
not outlying in the original space cannot be masked, which is the
conditioning on original-space outlyingness the proposal asks for. Labels
only say what a failure costs: a masked anomaly is a detection lost to the
projection.

\subsection*{Checks with known answers}

All 44 checks in \texttt{src/validate.py} pass. Among them: an identity
projection and a rotation move nothing; swapping the two spaces negates
$\lambda$ and swaps masked with phantom; a monotone rescoring changes
nothing; shuffled projected scores give an AO of 0 on average; comparing two
fresh re-measurements of the original data calls 2 to 3\% of crossings at
$\alpha = 5\%$; and of 20 outliers planted in the directions PCA discards,
17 to 20 are called masked over five seeds, while all 20 planted in the
plane PCA keeps are retained.

\subsection*{Results in two dimensions, top 5\%}

Masking rate MR, with the adjusted overlap AO in brackets. Means over
seeds: PCA is deterministic, the Gaussian random projection has 5
seeds, and UMAP, t-SNE and PaCMAP 3 each, at library defaults.
\begin{center}
\small
\begin{tabular}{lccccc}
\hline
 & shuttle & wbc & ionosphere & fault & magic \\
\hline
PCA & 0.57 ($+0.34$) & 0.00 ($+0.91$) & 1.00 ($-0.05$) & 0.58 ($+0.37$) & 0.39 ($+0.56$) \\
Random & 0.40 ($+0.49$) & 0.37 ($+0.58$) & 0.66 ($+0.31$) & 0.59 ($+0.36$) & 0.38 ($+0.58$) \\
UMAP & 0.91 ($+0.02$) & 0.61 ($+0.35$) & 1.00 ($-0.05$) & 0.82 ($+0.11$) & 0.85 ($+0.10$) \\
t-SNE & 0.73 ($+0.18$) & 0.81 ($+0.15$) & 0.85 ($+0.10$) & 0.70 ($+0.25$) & 0.93 ($+0.02$) \\
PaCMAP & 0.85 ($+0.07$) & 0.97 ($-0.03$) & 1.00 ($-0.05$) & 0.98 ($-0.03$) & 0.89 ($+0.06$) \\
\hline
\end{tabular}
\end{center}

\begin{itemize}
  \item UMAP, t-SNE and PaCMAP hide 61 to 100\% of the top 5\%, more than PCA
    or a random projection on four of the five datasets. PaCMAP's shortlist
    is close to chance on all five (AO at most 0.09).
  \item PCA hides from 0\% (wbc) to 100\% (ionosphere). The failures fall as
    PCA keeps more dimensions and vanish by full dimension on all five.
  \item A random projection hides less than PCA on shuttle and ionosphere,
    because PCA discards low-variance directions whatever they hold.
  \item Aggregate correlations miss the shortlist. On ionosphere PCA keeps a
    Spearman correlation of 0.56, and 0.55 for the weighted correlation of
    Schubert et al.\ (2012), while its 2D shortlist shares nothing with the
    original one.
  \item Shuttle is not a control for PCA, which hides 57\% of its top 5\%:
    its median LID is about 2 locally, but it spans six linear dimensions.
    It is also lattice data: all nine features are integers, and its 2D
    pictures under PCA, UMAP, t-SNE and PaCMAP each break into
    1,188 to 1,511 small islands of near-identical points
    (169 under a random projection), against 5 to 65 for magic.
  \item The neighbour-based methods gather isolated outliers into tight
    clumps. On wbc and fault their verdict depends on the neighbourhood size
    used to read the picture; on shuttle, ionosphere and magic it does not.
\end{itemize}

\subsection*{Open, for the supervisor}

\begin{enumerate}
  \item A resampling null per point, with permutation kept only for the
    chance level of AO, because the scDEED-style permutation null calls
    almost nothing. Is that acceptable?
  \item Budget: the top 5\% of points, which needs no labels, instead of the
    number of labelled anomalies?
  \item How to read the projection: the same $k$ as the original space (the
    proposal's rule), a profile over several $k$, or a multi-scale score?
  \item Keep shuttle as the case where local and global dimension differ?
  \item Lattice data: keep shuttle and wbc as they are and report the effect?
    The discrete fields of network data will raise the same question.
  \item Ranks, as now, or normalised scores (Kriegel et al.\ 2011) for
    version 2?
\end{enumerate}
.
% Numbers: results/week5_revision.json, results/week5_redundancy_v1.json,
% results/week5_score_check.json, results/week6_benchmark.json and
% results/week6_stats.json, written by src/revise_measure.py,
% src/score_check.py and src/benchmark.py.

\section*{Week 5: sanity checks and a multi-scale candidate}

All on the exploration set, at 2D with the top 5\% budget, with five seeds
for every stochastic method and TriMap added as a sixth method.

\subsection*{The candidate}

Version 1 reads each space with one neighbourhood size, $k = 20$. UMAP and
PaCMAP gather isolated outliers into tight clumps, and a member of a clump
of 40 has its 20th neighbour inside the clump: read at $k = 20$ the clump
looks dense, read at $k = 50$ it is seen to be apart. The candidate reads
both spaces at $K = \{20, 50, 100\}$,
\begin{equation*}
  s_i = \max_{k \in K} \frac{d_k(i)}{\operatorname{median}_j d_k(j)},
\end{equation*}
where $d_k(i)$ is the distance from $x_i$ to its $k$-th nearest neighbour,
and keeps everything else of version 1, which is the case $K = \{20\}$. The
precedent is Breunig, Kriegel, Ng and Sander (2000), who rank points by their
largest LOF over a range of MinPts.

\subsection*{Decision rule, fixed before the results were read}

Adopt the candidate if (i) it passes the checks with known answers, (ii) it
moves the masking rate of the linear methods by at most 0.10, (iii) its
verdict depends on its scales less than version 1's depends on $k$, and (iv)
the Jaccard overlap of its masked sets across seeds is at most 0.05 below
version 1's. Otherwise keep version 1.

\subsection*{Outcome: version 1 stays}

\begin{itemize}
  \item (i) holds. All 56 checks in \texttt{src/validate.py} pass. In a new
    known-answer test, forty outliers are placed in a picture as one tight
    clump far from the bulk: version 1 calls all forty hidden, the candidate
    none. Both call outliers placed inside the bulk hidden.
  \item (ii) fails narrowly: the other nine linear runs move by 0.083 at most,
    but shuttle under PCA moves from 0.569 to 0.468.
  \item (iii) holds: under three other sets of scales the candidate's
    masking rate stays within 0.10 in 89\% of runs (median change 0.010);
    version 1 under other $k$ does so in 71\% (median 0.041).
  \item (iv) fails narrowly: mean Jaccard 0.741 against 0.797. The loss is
    in t-SNE and TriMap.
\end{itemize}
The candidate is run beside version 1 on the benchmark, as a secondary
analysis declared before the run.

\subsection*{Version 1 against the candidate}

Masking rate, mean over five seeds; version 1, then the candidate.
\begin{center}
\small
\begin{tabular}{lccccc}
\hline
 & shuttle & wbc & ionosphere & fault & magic \\
\hline
PCA & 0.57, 0.47 & 0.00, 0.08 & 1.00, 1.00 & 0.58, 0.61 & 0.39, 0.37 \\
Random & 0.40, 0.38 & 0.37, 0.38 & 0.66, 0.66 & 0.59, 0.56 & 0.38, 0.38 \\
UMAP & 0.90, 0.90 & 0.77, 0.32 & 0.98, 1.00 & 0.81, 0.71 & 0.85, 0.71 \\
t-SNE & 0.73, 0.81 & 0.77, 0.27 & 0.84, 0.86 & 0.69, 0.77 & 0.93, 0.86 \\
PaCMAP & 0.85, 0.82 & 0.95, 0.17 & 1.00, 1.00 & 0.97, 0.62 & 0.89, 0.87 \\
TriMap & 0.61, 0.63 & 0.82, 0.22 & 0.97, 0.99 & 0.75, 0.72 & 0.94, 0.94 \\
\hline
\end{tabular}
\end{center}

\subsection*{Other findings}

\begin{itemize}
  \item \textbf{The verdicts depend on the detector.} LOF goes through the
    definition unchanged and stays close to calibrated, but orders the six
    methods differently from kNN on the three continuous datasets (rank
    correlations $+0.85$, $-0.43$ and $+0.26$). On magic, kNN finds the
    linear methods keep most of the top 5\% and the neighbour-based methods
    lose most of it; LOF finds every method loses 68 to 79\%.
  \item \textbf{Not redundant with continuity, except for linear maps.}
    For PCA and random projections, per-point continuity (Venna and Kaski)
    separates masked from retained points well (median AUC 0.90). For the
    neighbour-based methods it points the wrong way (median 0.24 under
    version 1): the points they hide keep their original neighbours better
    than the points they keep, as a group of outliers moved into the
    picture together and packed into a clump would.
  \item \textbf{Ranks, not scores.} Scores do not carry over between spaces
    of different dimension: under PCA the median score ratio of wbc's top
    5\% grows from 5.1 at nine dimensions to 20.1 at two, and Gaussian
    scaling (Kriegel, Kr\"oger, Schubert and Zimek 2011) saturates at the top.
  \item \textbf{Seeds.} Only 16\% of the top 5\% gets the same call from five
    random projections; one random projection says little about any single
    point. t-SNE is the steadiest method.
\end{itemize}

\section*{Week 6: the benchmark in two dimensions}

The fifteen benchmark datasets, capped at 20,000 points each, projected to
2D by PCA, a Gaussian random projection, UMAP, t-SNE, PaCMAP and TriMap
(library defaults, three seeds). The projections were computed while week 5
ran; no outlier score touched benchmark data until week 5 had decided.
Version 1 is the primary analysis, the multi-scale candidate a secondary one
declared before the run. The unit of comparison is the dataset, with one
value per method: the mean over seeds.

\subsection*{Masking rate, version 1, top 5\%}

\begin{center}
\small
\begin{tabular}{lcccccc}
\hline
 & PCA & Random & UMAP & t-SNE & PaCMAP & TriMap \\
\hline
cover & 0.58 & 0.66 & 0.85 & 0.79 & 0.82 & 0.77 \\
pageblocks & 0.34 & 0.31 & 0.85 & 0.87 & 0.84 & 0.55 \\
vowels & 0.79 & 0.79 & 0.72 & 0.54 & 0.82 & 0.75 \\
pendigits & 0.68 & 0.70 & 0.91 & 0.81 & 0.88 & 0.71 \\
cardio & 0.49 & 0.46 & 0.83 & 0.78 & 0.99 & 0.93 \\
waveform & 0.72 & 0.80 & 0.85 & 0.78 & 0.69 & 0.72 \\
fraud & 0.43 & 0.54 & 0.89 & 0.87 & 0.89 & 0.93 \\
wdbc & 0.05 & 0.37 & 0.60 & 0.58 & 0.44 & 0.49 \\
letter & 0.68 & 0.78 & 0.84 & 0.53 & 0.87 & 0.86 \\
satimage2 & 0.56 & 0.54 & 0.70 & 0.57 & 0.68 & 0.60 \\
celeba & 0.64 & 0.81 & 0.76 & 0.76 & 0.75 & 0.65 \\
campaign & 0.57 & 0.67 & 0.84 & 0.75 & 0.91 & 0.83 \\
mnist & 0.83 & 0.62 & 0.78 & 0.77 & 0.81 & 0.41 \\
musk & 0.58 & 0.70 & 0.64 & 0.52 & 0.54 & 0.47 \\
backdoor & 0.54 & 0.33 & 0.58 & 0.62 & 0.81 & 0.64 \\
\hline
mean & 0.57 & 0.60 & 0.78 & 0.70 & 0.78 & 0.69 \\
average rank & 2.27 & 2.67 & 4.73 & 3.33 & 4.60 & 3.40 \\
\hline
\end{tabular}
\end{center}

\subsection*{Comparison}

Friedman's test rejects equal masking rates ($\chi^2 = 21.4$, $p = 0.0007$;
Iman-Davenport $F = 5.6$). With six methods and fifteen datasets Nemenyi's
critical difference is 1.95 ranks. Pairwise Wilcoxon signed-rank tests with
Holm's correction, as Benavoli, Corani and Mangili (2016) recommend, find
that PCA hides less than UMAP (on 13 of 15 datasets) and PaCMAP (12 of 15),
the random projection less than UMAP (12 of 15), TriMap less than PaCMAP and
t-SNE less than UMAP. t-SNE and TriMap cannot be told from the linear
methods. The phantom rate and the adjusted overlap give the same groups.

\subsection*{Findings}

\begin{itemize}
  \item Every method hides most of the top 5\% on most datasets. PCA and
    the random projection hide more than half on 11 of 15, UMAP and t-SNE
    on all 15. No method is best everywhere: PCA hides least on six
    datasets, the random projection on four, t-SNE and TriMap on two each.
  \item Read at $k = 20$, the UMAP, t-SNE and PaCMAP pictures rank the
    labelled anomalies \emph{below} the normal points (AUC under 0.5) on
    9, 8 and 7 of the 15 datasets; PCA, the random projection and TriMap
    never do. Under the multi-scale reading this happens on 1, 0 and 2.
  \item Of the 1,496 labelled anomalies in the original top 5\%, the
    picture hides 50\% under PCA, 47\% under a random projection, 65\%
    under TriMap, 73\% under t-SNE, 81\% under UMAP and 84\% under PaCMAP.
  \item The phantom rate stays within 0.05 of the masking rate in 90\% of
    runs, as expected at 2D.
  \item The multi-scale candidate lowers the neighbour-based methods'
    masking rates by 0.09 to 0.11 on average and moves the linear ones by a
    median of 0.012. Under it t-SNE and TriMap move level with PCA, and only
    TriMap against PaCMAP stays significant. How a picture is read decides
    part of the answer, so which reading to keep matters.
\end{itemize}

\subsection*{Open, for the supervisor}

The full list, each question with the decision taken so the work could
continue, is in \texttt{docs/QUESTIONS.md}. The most pressing:
\begin{enumerate}
  \item Keep version 1, as the rule fixed in advance requires, or adopt the
    multi-scale candidate, which missed two of four criteria narrowly but
    reads clumps the way a person probably would?
  \item The verdicts depend on the detector: LOF orders the methods
    differently from kNN. Is kNN acceptable as the benchmark's detector,
    with LOF added on the continuous datasets in week 7?
  \item The benchmark split is still marked as proposed. Confirm it?
\end{enumerate}
 in main.tex after
% % docs/week4.tex
% Week 4: the outlier-faithfulness measure, version 1, tested on the
% exploration set. The same content as docs/MEASURE.md (sections 2, 4, 5
% and 7), shortened. Needs amsmath and amssymb only. There is no
% bibliography: references are written out in the text.
%
% Use: paste into progress.tex, or put % docs/week4.tex
% Week 4: the outlier-faithfulness measure, version 1, tested on the
% exploration set. The same content as docs/MEASURE.md (sections 2, 4, 5
% and 7), shortened. Needs amsmath and amssymb only. There is no
% bibliography: references are written out in the text.
%
% Use: paste into progress.tex, or put % docs/week4.tex
% Week 4: the outlier-faithfulness measure, version 1, tested on the
% exploration set. The same content as docs/MEASURE.md (sections 2, 4, 5
% and 7), shortened. Needs amsmath and amssymb only. There is no
% bibliography: references are written out in the text.
%
% Use: paste into progress.tex, or put \input{week4} in main.tex.
% Numbers: results/week4_measure.json and results/week4_lattice.json,
% written by src/explore_measure.py and src/lattice_check.py.

\section*{Week 4: outlier faithfulness, version 1}

Built and tested on the exploration set only: shuttle, wbc, ionosphere,
fault and magic. The repository calls these five the development set. No
benchmark dataset was used. Version 1 will change in week 5 and is frozen
in week 12.

\subsection*{Definition}

Let $x_1, \dots, x_n \in \mathbb{R}^D$ be the data and
$y_1, \dots, y_n \in \mathbb{R}^d$ a projection of it. The same detector
scores both spaces: $a_i$ is the distance from $x_i$ to its $k$-th nearest
neighbour among the other points, and $b_i$ the same for $y_i$, with
$k = 20$. Outlying in the projection therefore means ranked high by the
detector run on the projected coordinates.

\paragraph{Ranks.} $\rho_a(i)$ is the descending rank of $a_i$, with
$\rho_a(i) = 1$ for the most outlying point. Scores within a relative
tolerance of $10^{-9}$ count as tied, and tied points share the average of
the positions they span. $\rho_b$ is defined the same way.

\paragraph{Displacement,} in bits:
\begin{equation*}
  \lambda_i = \log_2 \frac{\rho_a(i)}{\rho_b(i)} .
\end{equation*}
A negative value means the projection moved the point towards the crowd, a
positive value towards the edge. $|\lambda_i|$ counts rank doublings, so a
move from rank 1 to rank 64 is $-6$ and a move from rank 1001 to rank 1064
is $-0.09$. The top of the ranking, where an analyst looks, dominates.

\paragraph{Shortlists.} With budget $m = \lceil 0.05\, n \rceil$,
$S_a = \{ i : \rho_a(i) \le m \}$ and $S_b = \{ i : \rho_b(i) \le m \}$. A
tied group shares one rank, so it is wholly in or wholly out.

\paragraph{Noise null.} For $r = 1, \dots, 100$, draw a random half $H_r$
of the points and score every point by its distance to the $(k/2)$-th
nearest neighbour within $H_r$, itself excluded. Ranking all $n$ points
gives $\rho^{(r)}$. The null for point $i$ is the set of differences
$D_i^{(rs)} = \log_2 \rho^{(r)}(i) - \log_2 \rho^{(s)}(i)$ over all pairs
$r < s$, and
\begin{equation*}
  p_i = \tfrac{1}{2} \Pr_{r<s}\bigl( |D_i^{(rs)}| \ge |\lambda_i| \bigr)
\end{equation*}
is the one-sided p-value in the direction the point moved. The null asks how
far the point's rank moves between two equally valid measurements of the
original data. It does not involve the projection, so it is computed once
per dataset.

\paragraph{Calls,} at level $\alpha = 0.05$:
\begin{center}
\begin{tabular}{lll}
\hline
category & rule & meaning \\
\hline
retained & $i \in S_a \cap S_b$ & kept on the shortlist \\
masked & $i \in S_a \setminus S_b$ and $p_i < \alpha$ & hidden, beyond noise \\
phantom & $i \in S_b \setminus S_a$ and $p_i < \alpha$ & promoted, beyond noise \\
ambiguous & crossed, $p_i \ge \alpha$ & crossed within noise \\
\hline
\end{tabular}
\end{center}

\paragraph{Aggregates.}
\begin{equation*}
  \mathrm{MR} = \frac{|M|}{|S_a|}, \qquad
  \mathrm{PR} = \frac{|P|}{|S_b|}, \qquad
  \mathrm{AO} = \frac{O - E}{\min(|S_a|, |S_b|) - E},
\end{equation*}
where $M$ and $P$ are the masked and phantom points, $O = |S_a \cap S_b|$,
and $E = |S_a|\,|S_b| / n$ is the overlap expected when the projected scores
are permuted across points. MR is the share of the detector's shortlist that
the projection hides, PR the share of the projection's shortlist that is not
outlying in the data. AO is 1 for identical shortlists and 0 at chance, the
same adjustment Campos et al.\ (2016) apply to precision at $n$.

\paragraph{Labels} are not used by the measure. A labelled anomaly that is
not outlying in the original space cannot be masked, which is the
conditioning on original-space outlyingness the proposal asks for. Labels
only say what a failure costs: a masked anomaly is a detection lost to the
projection.

\subsection*{Checks with known answers}

All 44 checks in \texttt{src/validate.py} pass. Among them: an identity
projection and a rotation move nothing; swapping the two spaces negates
$\lambda$ and swaps masked with phantom; a monotone rescoring changes
nothing; shuffled projected scores give an AO of 0 on average; comparing two
fresh re-measurements of the original data calls 2 to 3\% of crossings at
$\alpha = 5\%$; and of 20 outliers planted in the directions PCA discards,
17 to 20 are called masked over five seeds, while all 20 planted in the
plane PCA keeps are retained.

\subsection*{Results in two dimensions, top 5\%}

Masking rate MR, with the adjusted overlap AO in brackets. Means over
seeds: PCA is deterministic, the Gaussian random projection has 5
seeds, and UMAP, t-SNE and PaCMAP 3 each, at library defaults.
\begin{center}
\small
\begin{tabular}{lccccc}
\hline
 & shuttle & wbc & ionosphere & fault & magic \\
\hline
PCA & 0.57 ($+0.34$) & 0.00 ($+0.91$) & 1.00 ($-0.05$) & 0.58 ($+0.37$) & 0.39 ($+0.56$) \\
Random & 0.40 ($+0.49$) & 0.37 ($+0.58$) & 0.66 ($+0.31$) & 0.59 ($+0.36$) & 0.38 ($+0.58$) \\
UMAP & 0.91 ($+0.02$) & 0.61 ($+0.35$) & 1.00 ($-0.05$) & 0.82 ($+0.11$) & 0.85 ($+0.10$) \\
t-SNE & 0.73 ($+0.18$) & 0.81 ($+0.15$) & 0.85 ($+0.10$) & 0.70 ($+0.25$) & 0.93 ($+0.02$) \\
PaCMAP & 0.85 ($+0.07$) & 0.97 ($-0.03$) & 1.00 ($-0.05$) & 0.98 ($-0.03$) & 0.89 ($+0.06$) \\
\hline
\end{tabular}
\end{center}

\begin{itemize}
  \item UMAP, t-SNE and PaCMAP hide 61 to 100\% of the top 5\%, more than PCA
    or a random projection on four of the five datasets. PaCMAP's shortlist
    is close to chance on all five (AO at most 0.09).
  \item PCA hides from 0\% (wbc) to 100\% (ionosphere). The failures fall as
    PCA keeps more dimensions and vanish by full dimension on all five.
  \item A random projection hides less than PCA on shuttle and ionosphere,
    because PCA discards low-variance directions whatever they hold.
  \item Aggregate correlations miss the shortlist. On ionosphere PCA keeps a
    Spearman correlation of 0.56, and 0.55 for the weighted correlation of
    Schubert et al.\ (2012), while its 2D shortlist shares nothing with the
    original one.
  \item Shuttle is not a control for PCA, which hides 57\% of its top 5\%:
    its median LID is about 2 locally, but it spans six linear dimensions.
    It is also lattice data: all nine features are integers, and its 2D
    pictures under PCA, UMAP, t-SNE and PaCMAP each break into
    1,188 to 1,511 small islands of near-identical points
    (169 under a random projection), against 5 to 65 for magic.
  \item The neighbour-based methods gather isolated outliers into tight
    clumps. On wbc and fault their verdict depends on the neighbourhood size
    used to read the picture; on shuttle, ionosphere and magic it does not.
\end{itemize}

\subsection*{Open, for the supervisor}

\begin{enumerate}
  \item A resampling null per point, with permutation kept only for the
    chance level of AO, because the scDEED-style permutation null calls
    almost nothing. Is that acceptable?
  \item Budget: the top 5\% of points, which needs no labels, instead of the
    number of labelled anomalies?
  \item How to read the projection: the same $k$ as the original space (the
    proposal's rule), a profile over several $k$, or a multi-scale score?
  \item Keep shuttle as the case where local and global dimension differ?
  \item Lattice data: keep shuttle and wbc as they are and report the effect?
    The discrete fields of network data will raise the same question.
  \item Ranks, as now, or normalised scores (Kriegel et al.\ 2011) for
    version 2?
\end{enumerate}
 in main.tex.
% Numbers: results/week4_measure.json and results/week4_lattice.json,
% written by src/explore_measure.py and src/lattice_check.py.

\section*{Week 4: outlier faithfulness, version 1}

Built and tested on the exploration set only: shuttle, wbc, ionosphere,
fault and magic. The repository calls these five the development set. No
benchmark dataset was used. Version 1 will change in week 5 and is frozen
in week 12.

\subsection*{Definition}

Let $x_1, \dots, x_n \in \mathbb{R}^D$ be the data and
$y_1, \dots, y_n \in \mathbb{R}^d$ a projection of it. The same detector
scores both spaces: $a_i$ is the distance from $x_i$ to its $k$-th nearest
neighbour among the other points, and $b_i$ the same for $y_i$, with
$k = 20$. Outlying in the projection therefore means ranked high by the
detector run on the projected coordinates.

\paragraph{Ranks.} $\rho_a(i)$ is the descending rank of $a_i$, with
$\rho_a(i) = 1$ for the most outlying point. Scores within a relative
tolerance of $10^{-9}$ count as tied, and tied points share the average of
the positions they span. $\rho_b$ is defined the same way.

\paragraph{Displacement,} in bits:
\begin{equation*}
  \lambda_i = \log_2 \frac{\rho_a(i)}{\rho_b(i)} .
\end{equation*}
A negative value means the projection moved the point towards the crowd, a
positive value towards the edge. $|\lambda_i|$ counts rank doublings, so a
move from rank 1 to rank 64 is $-6$ and a move from rank 1001 to rank 1064
is $-0.09$. The top of the ranking, where an analyst looks, dominates.

\paragraph{Shortlists.} With budget $m = \lceil 0.05\, n \rceil$,
$S_a = \{ i : \rho_a(i) \le m \}$ and $S_b = \{ i : \rho_b(i) \le m \}$. A
tied group shares one rank, so it is wholly in or wholly out.

\paragraph{Noise null.} For $r = 1, \dots, 100$, draw a random half $H_r$
of the points and score every point by its distance to the $(k/2)$-th
nearest neighbour within $H_r$, itself excluded. Ranking all $n$ points
gives $\rho^{(r)}$. The null for point $i$ is the set of differences
$D_i^{(rs)} = \log_2 \rho^{(r)}(i) - \log_2 \rho^{(s)}(i)$ over all pairs
$r < s$, and
\begin{equation*}
  p_i = \tfrac{1}{2} \Pr_{r<s}\bigl( |D_i^{(rs)}| \ge |\lambda_i| \bigr)
\end{equation*}
is the one-sided p-value in the direction the point moved. The null asks how
far the point's rank moves between two equally valid measurements of the
original data. It does not involve the projection, so it is computed once
per dataset.

\paragraph{Calls,} at level $\alpha = 0.05$:
\begin{center}
\begin{tabular}{lll}
\hline
category & rule & meaning \\
\hline
retained & $i \in S_a \cap S_b$ & kept on the shortlist \\
masked & $i \in S_a \setminus S_b$ and $p_i < \alpha$ & hidden, beyond noise \\
phantom & $i \in S_b \setminus S_a$ and $p_i < \alpha$ & promoted, beyond noise \\
ambiguous & crossed, $p_i \ge \alpha$ & crossed within noise \\
\hline
\end{tabular}
\end{center}

\paragraph{Aggregates.}
\begin{equation*}
  \mathrm{MR} = \frac{|M|}{|S_a|}, \qquad
  \mathrm{PR} = \frac{|P|}{|S_b|}, \qquad
  \mathrm{AO} = \frac{O - E}{\min(|S_a|, |S_b|) - E},
\end{equation*}
where $M$ and $P$ are the masked and phantom points, $O = |S_a \cap S_b|$,
and $E = |S_a|\,|S_b| / n$ is the overlap expected when the projected scores
are permuted across points. MR is the share of the detector's shortlist that
the projection hides, PR the share of the projection's shortlist that is not
outlying in the data. AO is 1 for identical shortlists and 0 at chance, the
same adjustment Campos et al.\ (2016) apply to precision at $n$.

\paragraph{Labels} are not used by the measure. A labelled anomaly that is
not outlying in the original space cannot be masked, which is the
conditioning on original-space outlyingness the proposal asks for. Labels
only say what a failure costs: a masked anomaly is a detection lost to the
projection.

\subsection*{Checks with known answers}

All 44 checks in \texttt{src/validate.py} pass. Among them: an identity
projection and a rotation move nothing; swapping the two spaces negates
$\lambda$ and swaps masked with phantom; a monotone rescoring changes
nothing; shuffled projected scores give an AO of 0 on average; comparing two
fresh re-measurements of the original data calls 2 to 3\% of crossings at
$\alpha = 5\%$; and of 20 outliers planted in the directions PCA discards,
17 to 20 are called masked over five seeds, while all 20 planted in the
plane PCA keeps are retained.

\subsection*{Results in two dimensions, top 5\%}

Masking rate MR, with the adjusted overlap AO in brackets. Means over
seeds: PCA is deterministic, the Gaussian random projection has 5
seeds, and UMAP, t-SNE and PaCMAP 3 each, at library defaults.
\begin{center}
\small
\begin{tabular}{lccccc}
\hline
 & shuttle & wbc & ionosphere & fault & magic \\
\hline
PCA & 0.57 ($+0.34$) & 0.00 ($+0.91$) & 1.00 ($-0.05$) & 0.58 ($+0.37$) & 0.39 ($+0.56$) \\
Random & 0.40 ($+0.49$) & 0.37 ($+0.58$) & 0.66 ($+0.31$) & 0.59 ($+0.36$) & 0.38 ($+0.58$) \\
UMAP & 0.91 ($+0.02$) & 0.61 ($+0.35$) & 1.00 ($-0.05$) & 0.82 ($+0.11$) & 0.85 ($+0.10$) \\
t-SNE & 0.73 ($+0.18$) & 0.81 ($+0.15$) & 0.85 ($+0.10$) & 0.70 ($+0.25$) & 0.93 ($+0.02$) \\
PaCMAP & 0.85 ($+0.07$) & 0.97 ($-0.03$) & 1.00 ($-0.05$) & 0.98 ($-0.03$) & 0.89 ($+0.06$) \\
\hline
\end{tabular}
\end{center}

\begin{itemize}
  \item UMAP, t-SNE and PaCMAP hide 61 to 100\% of the top 5\%, more than PCA
    or a random projection on four of the five datasets. PaCMAP's shortlist
    is close to chance on all five (AO at most 0.09).
  \item PCA hides from 0\% (wbc) to 100\% (ionosphere). The failures fall as
    PCA keeps more dimensions and vanish by full dimension on all five.
  \item A random projection hides less than PCA on shuttle and ionosphere,
    because PCA discards low-variance directions whatever they hold.
  \item Aggregate correlations miss the shortlist. On ionosphere PCA keeps a
    Spearman correlation of 0.56, and 0.55 for the weighted correlation of
    Schubert et al.\ (2012), while its 2D shortlist shares nothing with the
    original one.
  \item Shuttle is not a control for PCA, which hides 57\% of its top 5\%:
    its median LID is about 2 locally, but it spans six linear dimensions.
    It is also lattice data: all nine features are integers, and its 2D
    pictures under PCA, UMAP, t-SNE and PaCMAP each break into
    1,188 to 1,511 small islands of near-identical points
    (169 under a random projection), against 5 to 65 for magic.
  \item The neighbour-based methods gather isolated outliers into tight
    clumps. On wbc and fault their verdict depends on the neighbourhood size
    used to read the picture; on shuttle, ionosphere and magic it does not.
\end{itemize}

\subsection*{Open, for the supervisor}

\begin{enumerate}
  \item A resampling null per point, with permutation kept only for the
    chance level of AO, because the scDEED-style permutation null calls
    almost nothing. Is that acceptable?
  \item Budget: the top 5\% of points, which needs no labels, instead of the
    number of labelled anomalies?
  \item How to read the projection: the same $k$ as the original space (the
    proposal's rule), a profile over several $k$, or a multi-scale score?
  \item Keep shuttle as the case where local and global dimension differ?
  \item Lattice data: keep shuttle and wbc as they are and report the effect?
    The discrete fields of network data will raise the same question.
  \item Ranks, as now, or normalised scores (Kriegel et al.\ 2011) for
    version 2?
\end{enumerate}
 in main.tex.
% Numbers: results/week4_measure.json and results/week4_lattice.json,
% written by src/explore_measure.py and src/lattice_check.py.

\section*{Week 4: outlier faithfulness, version 1}

Built and tested on the exploration set only: shuttle, wbc, ionosphere,
fault and magic. The repository calls these five the development set. No
benchmark dataset was used. Version 1 will change in week 5 and is frozen
in week 12.

\subsection*{Definition}

Let $x_1, \dots, x_n \in \mathbb{R}^D$ be the data and
$y_1, \dots, y_n \in \mathbb{R}^d$ a projection of it. The same detector
scores both spaces: $a_i$ is the distance from $x_i$ to its $k$-th nearest
neighbour among the other points, and $b_i$ the same for $y_i$, with
$k = 20$. Outlying in the projection therefore means ranked high by the
detector run on the projected coordinates.

\paragraph{Ranks.} $\rho_a(i)$ is the descending rank of $a_i$, with
$\rho_a(i) = 1$ for the most outlying point. Scores within a relative
tolerance of $10^{-9}$ count as tied, and tied points share the average of
the positions they span. $\rho_b$ is defined the same way.

\paragraph{Displacement,} in bits:
\begin{equation*}
  \lambda_i = \log_2 \frac{\rho_a(i)}{\rho_b(i)} .
\end{equation*}
A negative value means the projection moved the point towards the crowd, a
positive value towards the edge. $|\lambda_i|$ counts rank doublings, so a
move from rank 1 to rank 64 is $-6$ and a move from rank 1001 to rank 1064
is $-0.09$. The top of the ranking, where an analyst looks, dominates.

\paragraph{Shortlists.} With budget $m = \lceil 0.05\, n \rceil$,
$S_a = \{ i : \rho_a(i) \le m \}$ and $S_b = \{ i : \rho_b(i) \le m \}$. A
tied group shares one rank, so it is wholly in or wholly out.

\paragraph{Noise null.} For $r = 1, \dots, 100$, draw a random half $H_r$
of the points and score every point by its distance to the $(k/2)$-th
nearest neighbour within $H_r$, itself excluded. Ranking all $n$ points
gives $\rho^{(r)}$. The null for point $i$ is the set of differences
$D_i^{(rs)} = \log_2 \rho^{(r)}(i) - \log_2 \rho^{(s)}(i)$ over all pairs
$r < s$, and
\begin{equation*}
  p_i = \tfrac{1}{2} \Pr_{r<s}\bigl( |D_i^{(rs)}| \ge |\lambda_i| \bigr)
\end{equation*}
is the one-sided p-value in the direction the point moved. The null asks how
far the point's rank moves between two equally valid measurements of the
original data. It does not involve the projection, so it is computed once
per dataset.

\paragraph{Calls,} at level $\alpha = 0.05$:
\begin{center}
\begin{tabular}{lll}
\hline
category & rule & meaning \\
\hline
retained & $i \in S_a \cap S_b$ & kept on the shortlist \\
masked & $i \in S_a \setminus S_b$ and $p_i < \alpha$ & hidden, beyond noise \\
phantom & $i \in S_b \setminus S_a$ and $p_i < \alpha$ & promoted, beyond noise \\
ambiguous & crossed, $p_i \ge \alpha$ & crossed within noise \\
\hline
\end{tabular}
\end{center}

\paragraph{Aggregates.}
\begin{equation*}
  \mathrm{MR} = \frac{|M|}{|S_a|}, \qquad
  \mathrm{PR} = \frac{|P|}{|S_b|}, \qquad
  \mathrm{AO} = \frac{O - E}{\min(|S_a|, |S_b|) - E},
\end{equation*}
where $M$ and $P$ are the masked and phantom points, $O = |S_a \cap S_b|$,
and $E = |S_a|\,|S_b| / n$ is the overlap expected when the projected scores
are permuted across points. MR is the share of the detector's shortlist that
the projection hides, PR the share of the projection's shortlist that is not
outlying in the data. AO is 1 for identical shortlists and 0 at chance, the
same adjustment Campos et al.\ (2016) apply to precision at $n$.

\paragraph{Labels} are not used by the measure. A labelled anomaly that is
not outlying in the original space cannot be masked, which is the
conditioning on original-space outlyingness the proposal asks for. Labels
only say what a failure costs: a masked anomaly is a detection lost to the
projection.

\subsection*{Checks with known answers}

All 44 checks in \texttt{src/validate.py} pass. Among them: an identity
projection and a rotation move nothing; swapping the two spaces negates
$\lambda$ and swaps masked with phantom; a monotone rescoring changes
nothing; shuffled projected scores give an AO of 0 on average; comparing two
fresh re-measurements of the original data calls 2 to 3\% of crossings at
$\alpha = 5\%$; and of 20 outliers planted in the directions PCA discards,
17 to 20 are called masked over five seeds, while all 20 planted in the
plane PCA keeps are retained.

\subsection*{Results in two dimensions, top 5\%}

Masking rate MR, with the adjusted overlap AO in brackets. Means over
seeds: PCA is deterministic, the Gaussian random projection has 5
seeds, and UMAP, t-SNE and PaCMAP 3 each, at library defaults.
\begin{center}
\small
\begin{tabular}{lccccc}
\hline
 & shuttle & wbc & ionosphere & fault & magic \\
\hline
PCA & 0.57 ($+0.34$) & 0.00 ($+0.91$) & 1.00 ($-0.05$) & 0.58 ($+0.37$) & 0.39 ($+0.56$) \\
Random & 0.40 ($+0.49$) & 0.37 ($+0.58$) & 0.66 ($+0.31$) & 0.59 ($+0.36$) & 0.38 ($+0.58$) \\
UMAP & 0.91 ($+0.02$) & 0.61 ($+0.35$) & 1.00 ($-0.05$) & 0.82 ($+0.11$) & 0.85 ($+0.10$) \\
t-SNE & 0.73 ($+0.18$) & 0.81 ($+0.15$) & 0.85 ($+0.10$) & 0.70 ($+0.25$) & 0.93 ($+0.02$) \\
PaCMAP & 0.85 ($+0.07$) & 0.97 ($-0.03$) & 1.00 ($-0.05$) & 0.98 ($-0.03$) & 0.89 ($+0.06$) \\
\hline
\end{tabular}
\end{center}

\begin{itemize}
  \item UMAP, t-SNE and PaCMAP hide 61 to 100\% of the top 5\%, more than PCA
    or a random projection on four of the five datasets. PaCMAP's shortlist
    is close to chance on all five (AO at most 0.09).
  \item PCA hides from 0\% (wbc) to 100\% (ionosphere). The failures fall as
    PCA keeps more dimensions and vanish by full dimension on all five.
  \item A random projection hides less than PCA on shuttle and ionosphere,
    because PCA discards low-variance directions whatever they hold.
  \item Aggregate correlations miss the shortlist. On ionosphere PCA keeps a
    Spearman correlation of 0.56, and 0.55 for the weighted correlation of
    Schubert et al.\ (2012), while its 2D shortlist shares nothing with the
    original one.
  \item Shuttle is not a control for PCA, which hides 57\% of its top 5\%:
    its median LID is about 2 locally, but it spans six linear dimensions.
    It is also lattice data: all nine features are integers, and its 2D
    pictures under PCA, UMAP, t-SNE and PaCMAP each break into
    1,188 to 1,511 small islands of near-identical points
    (169 under a random projection), against 5 to 65 for magic.
  \item The neighbour-based methods gather isolated outliers into tight
    clumps. On wbc and fault their verdict depends on the neighbourhood size
    used to read the picture; on shuttle, ionosphere and magic it does not.
\end{itemize}

\subsection*{Open, for the supervisor}

\begin{enumerate}
  \item A resampling null per point, with permutation kept only for the
    chance level of AO, because the scDEED-style permutation null calls
    almost nothing. Is that acceptable?
  \item Budget: the top 5\% of points, which needs no labels, instead of the
    number of labelled anomalies?
  \item How to read the projection: the same $k$ as the original space (the
    proposal's rule), a profile over several $k$, or a multi-scale score?
  \item Keep shuttle as the case where local and global dimension differ?
  \item Lattice data: keep shuttle and wbc as they are and report the effect?
    The discrete fields of network data will raise the same question.
  \item Ranks, as now, or normalised scores (Kriegel et al.\ 2011) for
    version 2?
\end{enumerate}
.
% Numbers: results/week5_revision.json, results/week5_redundancy_v1.json,
% results/week5_score_check.json, results/week6_benchmark.json and
% results/week6_stats.json, written by src/revise_measure.py,
% src/score_check.py and src/benchmark.py.

\section*{Week 5: sanity checks and a multi-scale candidate}

All on the exploration set, at 2D with the top 5\% budget, with five seeds
for every stochastic method and TriMap added as a sixth method.

\subsection*{The candidate}

Version 1 reads each space with one neighbourhood size, $k = 20$. UMAP and
PaCMAP gather isolated outliers into tight clumps, and a member of a clump
of 40 has its 20th neighbour inside the clump: read at $k = 20$ the clump
looks dense, read at $k = 50$ it is seen to be apart. The candidate reads
both spaces at $K = \{20, 50, 100\}$,
\begin{equation*}
  s_i = \max_{k \in K} \frac{d_k(i)}{\operatorname{median}_j d_k(j)},
\end{equation*}
where $d_k(i)$ is the distance from $x_i$ to its $k$-th nearest neighbour,
and keeps everything else of version 1, which is the case $K = \{20\}$. The
precedent is Breunig, Kriegel, Ng and Sander (2000), who rank points by their
largest LOF over a range of MinPts.

\subsection*{Decision rule, fixed before the results were read}

Adopt the candidate if (i) it passes the checks with known answers, (ii) it
moves the masking rate of the linear methods by at most 0.10, (iii) its
verdict depends on its scales less than version 1's depends on $k$, and (iv)
the Jaccard overlap of its masked sets across seeds is at most 0.05 below
version 1's. Otherwise keep version 1.

\subsection*{Outcome: version 1 stays}

\begin{itemize}
  \item (i) holds. All 56 checks in \texttt{src/validate.py} pass. In a new
    known-answer test, forty outliers are placed in a picture as one tight
    clump far from the bulk: version 1 calls all forty hidden, the candidate
    none. Both call outliers placed inside the bulk hidden.
  \item (ii) fails narrowly: the other nine linear runs move by 0.083 at most,
    but shuttle under PCA moves from 0.569 to 0.468.
  \item (iii) holds: under three other sets of scales the candidate's
    masking rate stays within 0.10 in 89\% of runs (median change 0.010);
    version 1 under other $k$ does so in 71\% (median 0.041).
  \item (iv) fails narrowly: mean Jaccard 0.741 against 0.797. The loss is
    in t-SNE and TriMap.
\end{itemize}
The candidate is run beside version 1 on the benchmark, as a secondary
analysis declared before the run.

\subsection*{Version 1 against the candidate}

Masking rate, mean over five seeds; version 1, then the candidate.
\begin{center}
\small
\begin{tabular}{lccccc}
\hline
 & shuttle & wbc & ionosphere & fault & magic \\
\hline
PCA & 0.57, 0.47 & 0.00, 0.08 & 1.00, 1.00 & 0.58, 0.61 & 0.39, 0.37 \\
Random & 0.40, 0.38 & 0.37, 0.38 & 0.66, 0.66 & 0.59, 0.56 & 0.38, 0.38 \\
UMAP & 0.90, 0.90 & 0.77, 0.32 & 0.98, 1.00 & 0.81, 0.71 & 0.85, 0.71 \\
t-SNE & 0.73, 0.81 & 0.77, 0.27 & 0.84, 0.86 & 0.69, 0.77 & 0.93, 0.86 \\
PaCMAP & 0.85, 0.82 & 0.95, 0.17 & 1.00, 1.00 & 0.97, 0.62 & 0.89, 0.87 \\
TriMap & 0.61, 0.63 & 0.82, 0.22 & 0.97, 0.99 & 0.75, 0.72 & 0.94, 0.94 \\
\hline
\end{tabular}
\end{center}

\subsection*{Other findings}

\begin{itemize}
  \item \textbf{The verdicts depend on the detector.} LOF goes through the
    definition unchanged and stays close to calibrated, but orders the six
    methods differently from kNN on the three continuous datasets (rank
    correlations $+0.85$, $-0.43$ and $+0.26$). On magic, kNN finds the
    linear methods keep most of the top 5\% and the neighbour-based methods
    lose most of it; LOF finds every method loses 68 to 79\%.
  \item \textbf{Not redundant with continuity, except for linear maps.}
    For PCA and random projections, per-point continuity (Venna and Kaski)
    separates masked from retained points well (median AUC 0.90). For the
    neighbour-based methods it points the wrong way (median 0.24 under
    version 1): the points they hide keep their original neighbours better
    than the points they keep, as a group of outliers moved into the
    picture together and packed into a clump would.
  \item \textbf{Ranks, not scores.} Scores do not carry over between spaces
    of different dimension: under PCA the median score ratio of wbc's top
    5\% grows from 5.1 at nine dimensions to 20.1 at two, and Gaussian
    scaling (Kriegel, Kr\"oger, Schubert and Zimek 2011) saturates at the top.
  \item \textbf{Seeds.} Only 16\% of the top 5\% gets the same call from five
    random projections; one random projection says little about any single
    point. t-SNE is the steadiest method.
\end{itemize}

\section*{Week 6: the benchmark in two dimensions}

The fifteen benchmark datasets, capped at 20,000 points each, projected to
2D by PCA, a Gaussian random projection, UMAP, t-SNE, PaCMAP and TriMap
(library defaults, three seeds). The projections were computed while week 5
ran; no outlier score touched benchmark data until week 5 had decided.
Version 1 is the primary analysis, the multi-scale candidate a secondary one
declared before the run. The unit of comparison is the dataset, with one
value per method: the mean over seeds.

\subsection*{Masking rate, version 1, top 5\%}

\begin{center}
\small
\begin{tabular}{lcccccc}
\hline
 & PCA & Random & UMAP & t-SNE & PaCMAP & TriMap \\
\hline
cover & 0.58 & 0.66 & 0.85 & 0.79 & 0.82 & 0.77 \\
pageblocks & 0.34 & 0.31 & 0.85 & 0.87 & 0.84 & 0.55 \\
vowels & 0.79 & 0.79 & 0.72 & 0.54 & 0.82 & 0.75 \\
pendigits & 0.68 & 0.70 & 0.91 & 0.81 & 0.88 & 0.71 \\
cardio & 0.49 & 0.46 & 0.83 & 0.78 & 0.99 & 0.93 \\
waveform & 0.72 & 0.80 & 0.85 & 0.78 & 0.69 & 0.72 \\
fraud & 0.43 & 0.54 & 0.89 & 0.87 & 0.89 & 0.93 \\
wdbc & 0.05 & 0.37 & 0.60 & 0.58 & 0.44 & 0.49 \\
letter & 0.68 & 0.78 & 0.84 & 0.53 & 0.87 & 0.86 \\
satimage2 & 0.56 & 0.54 & 0.70 & 0.57 & 0.68 & 0.60 \\
celeba & 0.64 & 0.81 & 0.76 & 0.76 & 0.75 & 0.65 \\
campaign & 0.57 & 0.67 & 0.84 & 0.75 & 0.91 & 0.83 \\
mnist & 0.83 & 0.62 & 0.78 & 0.77 & 0.81 & 0.41 \\
musk & 0.58 & 0.70 & 0.64 & 0.52 & 0.54 & 0.47 \\
backdoor & 0.54 & 0.33 & 0.58 & 0.62 & 0.81 & 0.64 \\
\hline
mean & 0.57 & 0.60 & 0.78 & 0.70 & 0.78 & 0.69 \\
average rank & 2.27 & 2.67 & 4.73 & 3.33 & 4.60 & 3.40 \\
\hline
\end{tabular}
\end{center}

\subsection*{Comparison}

Friedman's test rejects equal masking rates ($\chi^2 = 21.4$, $p = 0.0007$;
Iman-Davenport $F = 5.6$). With six methods and fifteen datasets Nemenyi's
critical difference is 1.95 ranks. Pairwise Wilcoxon signed-rank tests with
Holm's correction, as Benavoli, Corani and Mangili (2016) recommend, find
that PCA hides less than UMAP (on 13 of 15 datasets) and PaCMAP (12 of 15),
the random projection less than UMAP (12 of 15), TriMap less than PaCMAP and
t-SNE less than UMAP. t-SNE and TriMap cannot be told from the linear
methods. The phantom rate and the adjusted overlap give the same groups.

\subsection*{Findings}

\begin{itemize}
  \item Every method hides most of the top 5\% on most datasets. PCA and
    the random projection hide more than half on 11 of 15, UMAP and t-SNE
    on all 15. No method is best everywhere: PCA hides least on six
    datasets, the random projection on four, t-SNE and TriMap on two each.
  \item Read at $k = 20$, the UMAP, t-SNE and PaCMAP pictures rank the
    labelled anomalies \emph{below} the normal points (AUC under 0.5) on
    9, 8 and 7 of the 15 datasets; PCA, the random projection and TriMap
    never do. Under the multi-scale reading this happens on 1, 0 and 2.
  \item Of the 1,496 labelled anomalies in the original top 5\%, the
    picture hides 50\% under PCA, 47\% under a random projection, 65\%
    under TriMap, 73\% under t-SNE, 81\% under UMAP and 84\% under PaCMAP.
  \item The phantom rate stays within 0.05 of the masking rate in 90\% of
    runs, as expected at 2D.
  \item The multi-scale candidate lowers the neighbour-based methods'
    masking rates by 0.09 to 0.11 on average and moves the linear ones by a
    median of 0.012. Under it t-SNE and TriMap move level with PCA, and only
    TriMap against PaCMAP stays significant. How a picture is read decides
    part of the answer, so which reading to keep matters.
\end{itemize}

\subsection*{Open, for the supervisor}

The full list, each question with the decision taken so the work could
continue, is in \texttt{docs/QUESTIONS.md}. The most pressing:
\begin{enumerate}
  \item Keep version 1, as the rule fixed in advance requires, or adopt the
    multi-scale candidate, which missed two of four criteria narrowly but
    reads clumps the way a person probably would?
  \item The verdicts depend on the detector: LOF orders the methods
    differently from kNN. Is kNN acceptable as the benchmark's detector,
    with LOF added on the continuous datasets in week 7?
  \item The benchmark split is still marked as proposed. Confirm it?
\end{enumerate}
